\section{Setup and Preliminaries}

\paragraph{Notations.} We use $\|\cdot\|$ to denote the Euclidean 2-norm. For a symmetric positive definite matrix $A$ and a vector $x$ of proper dimension, we let $\|x\|_{A}\coloneqq \sqrt{x^\top A x}$ be the weighted 2-norm (or the $A$-norm). We let $\dotp{\cdot}{\cdot}$ denote the inner product in Euclidean  space such that $\dotp{x}{y} = x^\top y$. We let $\cB_d $ denote the unit ball in $\RR^d$, and $\cS_{d-1} = \{x\in\RR^d:\|x\| \leq  1\}$ denote the unit sphere in  $\RR^d$. For an interger $N\geq 1$, we let $[N]$ denote the set $\{1,2,\dots,N\}$. We use the $\cO(\cdot	)$ notation to suppress problem-dependent constants, and the $\tilde{\cO}(\cdot)$ notation further suppresses polylog factors.

\paragraph{Problem Formulation.} We consider the stochastic linear bandit problem. Let $\theta^\star \in \RR^d$ be a fixed but unknown parameter. At each time $t\in[T]$, a policy $\pi$  selects action $x_t$ from a set of action  $\cX_t \subset \RR^d $ according to the past observations and receives a reward $r_t  = \dotp{x_t}{\theta^\star} + \varepsilon_t$, where $ \varepsilon_t$ is mean-zero noise with a distribution specified in Assumption \ref{assump:noise} below. We measure the performance of $\pi$ with the cumulative expected regret 
%
$\cR(T) = \sum_{t= 1}^T\dotp{x^{\star}_t}{\theta^\star} -  \dotp{x_t}{\theta^\star},$
%
where $x^{\star}_t$ is the best action at time $t$, i.e., $x^{\star}_t=\arg\max_{x \in \cX_t}\dotp{x}{\theta^\star}$.
Let $\mathcal{F}_t$ be a $\sigma$-algebra generated by the history $(x_1, r_1, \ldots, x_t, r_t)$ and the prior knowledge, $\mathcal{F}_0$. Therefore, $\{\mathcal{F}_t\}_{t\ge 0}$ forms a filteration such that each $\mathcal{F}_t$ encodes all the information up to the end of period $t$.



\subsection{Assumptions}
We make the following assumptions that are standard in the relevant literature.

\begin{assumption}[Bounded parameter]\label{assump:bounded-parameter}
	The unknown parameter  $\param$ is bounded as  $\|\param\|\leq S$, where $S > 0$ is known.
\end{assumption}
\begin{assumption}[Bounded action sets]\label{assump:bounded-action-sets}
	The action sets $\{\cX_t\}$ are uniformly bounded and closed subsets of $\RR^d$, such that $\|x\|\leq X_t$ for all $x\in \cX_t$ and all $t\in[T]$, where $X_t$'s are known and $\sup_{t\geq 1}\left\{X_t\right\} < \infty $. Without loss of generality, we normalize the actions such that $\cX_t \subset \cS_{d-1}$ for all $t\in[T]$.
\end{assumption}


\begin{assumption} [Noise]\label{assump:noise}
	The noise sequence $\{\varepsilon_t\}_{t\ge 1}$ is conditionally mean-zero and $R$-subgaussian, where $R$ is known. Formally, 
	for all real valued $\lambda$, $ \EE\left[e^{\lambda \varepsilon_t} | \cF_t\right] \leq \exp \left(\lambda^2 R^2 / 2\right)$. 	This condition implies that $\EE \left[\varepsilon_t  | \cF_t\right]=0$ for all $t\geq 1$.
\end{assumption}



\subsection{Regularized  Least Square and Confidence Ellipsoid}

Here, we review useful frequentist tools developed in \cite{abbasi2011improved}.

Let $(x_1,\dots,x_t)$ be an arbitrary sequence of actions and $(r_1,\dots,r_t)$ be the corresponding rewards. In LB problems, the parameter $\theta^*$ is estimated with the regularized ordinary least-square estimator. Fix a regularization parameter $\lambda$, the design matrix, and the Regularized  Least Square (RLS) estimate are respectively given by
\#\label{eq:ridge}
V_t=\lambda I +\sum_{s=1}^{t} x_s x_s^{\top}, \quad \hat{\theta}_t=V_t^{-1} \sum_{s=1}^{t} x_s r_{s}.
\#

The following proposition in \cite{abbasi2011improved} states that the RLS estimate $\hparam_t$  concentrates around $\param$ properly.

\begin{proposition}[Theorem 2 in \cite{abbasi2011improved}]
	\label{prop:beta}
	Fix $\delta  \in (0,1)$, it holds with probability $1-\delta$ that
	\$
	&\left\|\widehat{\theta}_t-\theta^{\star}\right\|_{V_t} \leq \beta_t^{RLS}(\delta), \quad \left|\dotp{x}{\widehat{\theta}_t-\theta^{\star}}\right| \leq\|x\|_{V_t^{-1}} \beta_t^{RLS}(\delta)
	\$
	for all $x\in \mathbb{R}^d$,   where
	\#
	\beta_t^{RLS}(\delta) = R \sqrt{2 \log \frac{(\lambda+t)^{d / 2} \lambda^{-d / 2}}{\delta}}+\sqrt{\lambda} S.
	\#
\end{proposition}
Hereafter, we might omit the dependence of $\beta_t$ on $\delta$ if there is no ambiguity. Propostion~\ref{prop:beta} enables us to construct the following sequence of confidence ellipsoids.

\begin{definition}
	\label{def:rls}
	Fix $\delta\in (0,1)$. We define the RLS confidence ellipsoid as	$
	\cE_t^{RLS}(\delta)  = \{\theta \in \RR^d: \|\theta - \hparam_t\|_{V_t}\leq \beta_t^{RLS}(\delta) \}$.

\end{definition}


% todo: add ellipse potential references
The next proposition, known as the \emph{elliptical potential lemma}, presents a concentration result that plays a central role in bounding the regret. Basically, this proposition is saying that the cumulated prediction error incurred by the action sequence used to estimate $\param$ is small.
\begin{proposition}[Lemma 11 in \cite{abbasi2011improved}] \label{prop:potential}
	If $\lambda >1$, for an arbitrary sequence $(x_1,\dots,x_t)$, it holds that
	$
	\sum_{s=1}^t\left\|x_s\right\|_{V_s^{-1}}^2 \leq 2 \log \frac{\operatorname{det}\left(V_{t+1}\right)}{\operatorname{det}(\lambda I)} \leq 2 d \log \left(1+\frac{t}{\lambda}\right).
	$

\end{proposition}


\section{POFUL Algorithms}
\label{sec:POFUL}
In this section, we introduce a generalized framework of  OFUL that we call POFUL (Pivot OFUL). POFUL takes as input a sequence of inflation parameters $\{\iota_t\}_{t\in [T]}$,  feasible (randomized) pivots $\{\tilde{\theta}_t\}_{t\in[T]}$ and optimism parameters $\{\tau_t\}_{t\in [T]}$. The inflation parameters are used to construct confidence ellipsoids that contain $\{\tilde{\theta}_t\}_{t\in[T]}$ with high probability. This turns into the following condition for $\{\iota_t\}_{t\in [T]}$ and $\{\tilde{\theta}_t\}_{t\in[T]}$.

\begin{definition} \label{def:tilde_beta}
Fix  $\delta \in (0,1)$ and  $\delta^\prime = \delta/2 T$. Given the inflation parameters $\{\iota_t\}_{t\in [T]}$, we call random variables $\{\tilde{\theta}_t\}_{t\in[T]}$  feasible pivots if for all $t\in[T]$,
\$
	\PP\left[\tilde{\theta}_t \in \cE_t^{\PVT}(\delta ^\prime) |\mathcal{F}_t\right] \geq 1 - \delta^\prime,
	\$ where we define the \emph{pivot ellipsoid} as
$
 \cE_t^{\PVT} (\delta) \coloneqq \left \{\theta \in \RR^d: \|\theta - \hparam_t\|_{V_t}\leq \iota_t \beta^{RLS}_t (\delta ) \right\}$.
\end{definition}

At each time $t$, POFUL chooses the action that maximizes the optimistic reward \#\label{eq:POFUL}
\tilde{x}_t = \argmax_{x\in\cX_t} \  \dotp{x}{\tilde{\theta}_{t}}  +   \tau_t  \|x\|_{V^{-1}_t}\beta_t^{RLS}(\delta^\prime).
\#

A pseudocode representation of POFUL is provided in Algorithm~\ref{alg:POFUL}.

\begin{algorithm}
	\caption{POFUL }\label{alg:POFUL}
	\begin{algorithmic}
		% \Require $n \geq 0$
		% \Ensure $y = x^n$
		\REQUIRE $T$, $\delta$,  $\{\iota_t\}_{t\in [T]}$, $\{\tau_t\}_{t\in [T]}$
		\STATE Initialize $ V_{0} \leftarrow \lambda \mathbb{I}$, $\widehat{\theta}_{1} \leftarrow 0$,  $\delta^\prime \leftarrow \delta/2T$
		\FOR{t = 1, 2 , \dots, T }
		\STATE Sample a feasible pivot $\tilde{\theta}_{t} $ with respect to  $\iota_t$ according to Definition~\ref{def:tilde_beta}
		% \STATE Sample $\tilde{\theta}_{t} \sim \mathcal{N}\left(\hat{\theta}_{t}, \iota^{2}_t  V_{t}^{-1}\right)$
		\STATE $\tilde{x}_{t} \leftarrow \argmax_{x\in\cX_t} \dotp{x}{\tilde{\theta}_{t}}  +   \tau_t  \|x\|_{V^{-1}_t}\beta_t^{RLS}(\delta^\prime)$
		\STATE Observe reward $r_{t}$
		\STATE $V_{t+1} \leftarrow V_{t}+\tilde{x}_{t}  \tilde{x}_{t}^{\top}$
		\STATE $\hat{\theta}_{t+1} \leftarrow V_{t+1}^{-1}\sum_{s=1}^{t} \tilde{x}_s r_{s}$.
		\ENDFOR
	\end{algorithmic}
\end{algorithm}






The intuition behind POFUL is combining the exploration mechanism of OFUL and LinTS. OFUL encourages exploration by introducing the uncertainty term $\tau_t  \|x\|_{V^{-1}_t}\beta_t^{RLS}(\delta^\prime)$ in the reward, while LinTS explores by sampling from the confidence ellipsoid. We allow POFUL to select an arbitrary pivot from $\cE_t^{\PVT}$ to encompass arbitrary exploration mechanisms within $\cE_t^{\PVT}$.  By definition, POFUL includes OFUL, LinTS, TS-Freq, and Greedy as special cases. 
\paragraph{OFUL} 
For stochastic linear bandit problems, OFUL chooses actions by solving the optimization problem 
	$
	\max_{x\in\cX_t}  \   \dotp{x}{ \hat{\theta}_t}  + \|x\|_{V_t^{-1}}   \beta^{RLS}_t(\delta^\prime).
	$
Therefore, OFUL is a specially case of POFUL where $\iota_t = 0$, $\tau_t = 1$ and $
\tilde{\theta}_t = \hat{\theta}_t$, the center of the confidence ellipsoid, for all $t\in[T]$.

\paragraph{LinTS and TS-Freq}	
Linear Thompson Sampling (LinTS) algorithm is a generic randomized algorithm that samples from a distribution constructed from the RLS estimate at each step. At time $t$, LinTS samples as $\tilde{\theta}_{t}  =\hat{\theta}_t +\iota^{TS}_t \beta_t^{RLS} (\delta^\prime) V_t^{-\frac{1}{2}} \eta_t$, where $\delta^\prime = \delta/2T$, and $\iota^{TS}_t$ is an inflation parameter controlling the scale of the sampling range, and $\eta_t$ is a random sample from a normalized sampling distribution $\cD^{SA}(\delta^\prime)$ that concentrates with high probability. 


\begin{definition}
	 Fix any $\delta \in (0,1)$. $\cD^{SA} (\delta) $  is a distribution such that $
	\PP_{\theta \sim \cD^{SA}(\delta)}\left[\|\theta - \hparam_t\|_{V_t} \leq 1\right] \geq 1 - \delta.
	$
	\end{definition}
LinTS is a special case of POFUL where  $\iota_t = \iota^{TS}_t$, $\tau_t = 0$ and $\tilde{\theta_t} = \hat{\theta}_t +\iota^{TS}_t \beta_t^{RLS} (\delta^\prime) V_t^{-\frac{1}{2}} \eta_t$. We remark that setting the inflation parameter $\iota_t = \tilde{\cO}(1)$ corresponds to the original LinTS algorithm. On the other hand,  setting $\iota_t= \tilde{\cO}(\sqrt{d})$  corresponds to the frequentist variant of LinTS studied in \cite{agrawal2013thompson,abeille2017linear}, denoted by TS-Freq. This means TS-Freq inflates the posterior by a factor of order $\sqrt{d}$.
\paragraph{Greedy} Greedy is a special case of POFUL where $\iota_t = \tau_t = 0$ and $\tilde{\theta}_t = \hat{\theta}_t$  for all $t\in[T]$.


\begin{comment}
	A pseudocode representation of LinTS is provided in Algorithm~\ref{alg:TS}.



Define $\beta_t^{TS}(\delta) = \iota_t \beta_t^{RLS}(\delta)$. By construction, we have the following confidence ellipsoid for the LinTS samples $\{\tparam_t\}$.


\begin{proposition}
	\label{prop:ts}
Fix $\delta\in (0,1)$ and $\delta^\prime = \delta/2T$. Define \$\beta^{TS}_t (\delta )\coloneqq   \iota_t \beta_t^{RLS}(\delta) \$ and the ellipsoid \#
	\cE_t^{TS} =\left \{\theta \in \RR^d: \|\theta - \hparam_t\|_{V_t}\leq \beta^{TS}_t (\delta^\prime ) \right\}.
	\# Consider the concentration event $\tilde{\cA}_t = \left\{ \forall s\leq t,   \tilde{\theta}_s \in \cE_s^{TS}\right\}$ for $t \in [1,T]$ . 
 It holds that $\prob{\tilde{\cA}_t}  \geq 1 - \frac{\delta}{2}$.
	\end{proposition}
\$
\dotp{x^{\star}_t}{ \theta^{\star}} - \dotp{\tilde{x}_t}{  \tilde{\theta}_t}&\leq \dotp{x^{\star}_t}{ \tilde{\theta}_t}+U^{RLS}_t(x^{\star}_t) + U^{\PVT}_t(x^{\star}_t) -  \dotp{\tilde{x}_t}{  \tilde{\theta}_t}\\
&\leq U^{RLS}_t(x^{\star}_t) + U^{\PVT}_t(x^{\star}_t) + U^{\PVT}_t(\tilde{x}_t) -  U^{\PVT}_t(x_t^{\star})\\
&\leq U^{RLS}_t(x^{\star}_t) +  U^{\PVT}_t(\tilde{x}_t) 
\$
\end{comment}






\begin{comment}
A pseudocode representation of LinTS is provided in Algorithm~\ref{alg:TS}, with a slight generalization that allows the posterior at time $t$ to be inflated by a factor $\iota_t$.
 For this purpose, we define $\cD^{TS}_t$ as below.





Specifically, at each time $t$, LinTS samples from a distribution $\cD^{TS}_t$ satisfying the following condition:

\begin{definition}
	 $\cD^{TS}_t = \cD^{TS}_t( \hparam_t, V_t, \beta_t^{TS})$  is a multivariate distribution that has radius $\beta_t^{TS}$ if for any $\delta\in(0,1)$ it holds that $
	\PP_{\theta \sim \cD^{TS}_t}\left[\theta \in \cE_t^{TS}(\delta)\right] \geq 1 - \delta,
	$
	where we define the ellipsoid \#
	\cE_t^{TS} (\delta)\coloneqq\left \{\theta \in \RR^d: \|\theta - \hparam_t\|_{V_t}\leq \beta^{TS}_t (\delta) \right\}.
	\#

	
	\end{definition}
Bayesian algorithm that assigns $\param$ a Gaussian prior, assumes noise has a Gaussian distribution, and updates the posterior according to the observations using Bayes rule. Suppose the prior of $\param$ at time $t$ is given by $\cN(\hparam_t, {\color{red} R^2}  V_t) $ where $\hparam_t, V_t$ is defined in \eqref{eq:ridge}, it follows from direct computation that the posterior at time $t+1$ is given by $\cN(\hparam_{t+1}, {\color{red} R^2} V_{t+1}) $, i.e.  the updates of $\hat{\theta}_{t}$ and $V_t$ coincide with \eqref{eq:ridge}. Hence TS can be implemented by sampling $\tparam_t$ from $\mathcal{N}\left(\hat{\theta}_{t}, {\color{red} R^2} V_{t}^{-1}\right)$ at each $t$ and then taking the corresponding optimal action ${x}_{t}=\arg \max _{x \in \mathcal{X}}\dotp{ x}{ \tilde{\theta}_{t}}$. 






We remark that setting the inflation parameter $\iota_t \equiv {\color{red} R \sqrt{2 d \log (\frac{2d}{ \delta })  }/\beta^{RLS}_{t}(\delta)}  = \tilde{\cO}(1)$ corresponds to the original LinTS algorithm. On the other hand,  setting $\iota_t = \sqrt{2 d \log (\frac{2d}{ \delta })  } = \tilde{\cO}(\sqrt{d})$  corresponds to the frequentist variant of LinTS studied in \cite{abeille2017linear}, denoted by TS-Freq. This means TS-Freq inflates the posterior by a factor of order $\sqrt{d}$.
  



\begin{algorithm}
	\caption{Linear Thompson Sampling }\label{alg:TS}
	\begin{algorithmic}
		% \Require $n \geq 0$
		% \Ensure $y = x^n$
		\REQUIRE $T$, $\delta$,  inflation  parameters $\{\iota_t\}_{t\in [T]}$ 
		\STATE Initialize $ V_{0} \leftarrow \lambda \mathbb{I}$, $\widehat{\theta}_{1} \leftarrow 0$ and set $\delta^\prime = \delta/2T$
		\FOR{t = 1, 2 , \dots, T }
		\STATE Sample $\eta_t\sim \cD^{SA}(\delta^\prime)$ 
		\STATE $\tilde{\theta}_{t}  \leftarrow \hat{\theta}_t +\iota_t \beta_t^{RLS} (\delta^\prime) V_t^{-\frac{1}{2}} \eta_t$
		% \STATE Sample $\tilde{\theta}_{t} \sim \mathcal{N}\left(\hat{\theta}_{t}, \iota^{2}_t  V_{t}^{-1}\right)$
		\STATE ${x}_{t} \leftarrow \arg \max _{x \in \mathcal{X}}\dotp{ x}{ \tilde{\theta}_{t}}$
		\STATE Observe reward $r_{t}$
		\STATE $V_{t+1} \leftarrow V_{t}+{x}_{t}  {x}_{t}^{\top}$
		\STATE $\hat{\theta}_{t+1} \leftarrow V_{t+1}^{-1}\sum_{s=1}^{t} x_s r_{s}$.
		\ENDFOR
	\end{algorithmic}
\end{algorithm}


\end{comment}

















% todo add a remark that this is equivalent to adding a random vector


